\chapter{Verificación independiente sin cuenta}\label{cap:verificacion-independiente}

\paraque Permite a cualquier persona (perito, juez, contraparte) comprobar un
paquete de evidencia de un documento sellado sin entrar a Qadra.

\disponible Cualquier persona con el archivo ZIP.

\textbf{Necesita:}
\begin{itemize}
  \item El ZIP de evidencia original, sin modificar.
  \item Una herramienta para descomprimir archivos ZIP.
  \item OpenSSL versión 3 instalado (referencia: 3.5.7). No se necesita
        conexión a internet.
\end{itemize}

\pasos\ Abra una terminal en la carpeta del ZIP y ejecute, en orden:
\begin{enumerate}
  \item \textbf{Descomprimir el paquete.}
\begin{verbatim}
unzip -o evidencia.zip -d evidencia
cd evidencia
\end{verbatim}
  \item \textbf{Integridad.} Calcule la huella del documento y compárela con
        la que indica el archivo \archivo{INSTRUCCIONES.md}. Deben ser
        idénticas.
\begin{verbatim}
openssl dgst -sha256 <documento>
\end{verbatim}
  \item \textbf{Firma.} Debe responder \texttt{Verified OK}.
\begin{verbatim}
openssl x509 -in certificado.pem -pubkey -noout > firmante.pub.pem
openssl dgst -sha256 -verify firmante.pub.pem \
  -signature <documento>.sig <documento>
\end{verbatim}
  \item \textbf{Certificado y revocación.} Debe responder
        \texttt{certificado.pem: OK}.
\begin{verbatim}
cat ca.pem crl.pem > ca-y-crl.pem
openssl verify -crl_check -CAfile ca-y-crl.pem certificado.pem
\end{verbatim}
  \item \textbf{Sello de tiempo.} Debe responder \texttt{Verification: OK}.
        Si no existe \archivo{tsa-chain.pem}, use \archivo{ca.pem}.
\begin{verbatim}
openssl ts -verify -data <documento> -in <documento>.tsr \
  -CAfile tsa-chain.pem
\end{verbatim}
\end{enumerate}

Sustituya \texttt{<documento>} por el nombre del archivo que viene en el
paquete.

\begin{figure}[H]
  \centering
  \capturapendiente{capturas/B-05-verificacion-openssl.png}%
    {Terminal con Verified OK, certificado.pem: OK y Verification: OK · sin cuenta}
  \caption{Verificación independiente con OpenSSL.}
  \label{fig:manual-verificacion-openssl}
\end{figure}

\resultado Las cuatro comprobaciones responden correctamente: el documento no
cambió, la firma corresponde al documento, el certificado era válido y el
sello de tiempo es auténtico.

\begin{itemize}
  \item Si el documento tiene firma personal, \archivo{INSTRUCCIONES.md}
        indica dos comprobaciones adicionales; sígalas tal como aparecen.
  \item La lista de revocación del paquete refleja el estado al momento de
        sellar. Para el estado actual, pida una lista vigente a la CA
        emisora.
  \item Confirme por un medio independiente que el certificado raíz
        (\archivo{ca.pem}) es el del despacho.
  \item La explicación detallada de cada paso está en la
        \emph{Guía de verificación independiente con OpenSSL} del reporte
        técnico del Trabajo Terminal.
\end{itemize}

\importante{El sello de tiempo de esta demostración proviene de una
autoridad local: no equivale a una constancia NOM-151 emitida por un PSC
autorizado.}
